\chapter{原子核}
\problem{10.1}{${}^1_1\mathrm H$ 和 ${}^1_0\mathrm n$ 的质量分别为 1.0078252 和 1.0086654 质量单位，算出 ${}^{12}_6\mathrm C$ 中每核子的平均结合能（1 原子质量单位 $=931.5\,\mathrm{MeV}/c^2$）。}
\answerarea
\problem{10.2}{从下列各种粒子的质量数据中选用需要的数值，算出 ${}^{30}_{14}\mathrm{Si}$ 中每核子的平均结合能：
\[
\begin{array}{c@{\quad}r@{\qquad}c@{\quad}r}
\mathrm e&0.000548&{}^2_1\mathrm H&2.014102\\
{}^1_0\mathrm n&1.008665&{}^{30}_{14}\mathrm{Si}&29.973786\\
{}^1_1\mathrm H&1.007825&&
\end{array}
\]}
\answerarea\clearpage
\problem{10.3}{${}^{232}_{90}\mathrm{Th}$ 放射 $\alpha$ 射线成为 ${}^{228}_{88}\mathrm{Ra}$。从含有 1 克 ${}^{232}_{90}\mathrm{Th}$ 的一片薄膜测得每秒放射 4100 粒 $\alpha$ 粒子，试算出 ${}^{232}_{90}\mathrm{Th}$ 的半衰期为 $1.4\times10^{10}$ 年。}
\answerarea
\problem{10.4}{在考古工作中，可以从古生物遗骸中 ${}^{14}\mathrm C$ 的含量推算古生物到现在的时间 $t$。设 $\rho$ 是古生物遗骸中 ${}^{14}\mathrm C$ 和 ${}^{12}\mathrm C$ 存量之比，$\rho_0$ 是空气中 ${}^{14}\mathrm C$ 和 ${}^{12}\mathrm C$ 存量之比，试推导出下列公式：
\[
t=\frac{T\ln(\rho_0/\rho)}{\ln2},
\]
式中 $T$ 为 ${}^{14}\mathrm C$ 的半衰期。}
\answerarea\clearpage
\problem{10.5}{核力在原子核大小的距离内是很强的吸引力，它克服了质子间的（元素氢除外，那里只有一粒质子）库仑推斥力的作用而使原子核结合着。足见在原子核中核力的作用超过质子间的库仑推斥力作用；从质子间推斥力的大小可以粗略地了解到核力大小的低限。试计算原子核中两粒质子之间的库仑推斥力的大小（用公斤表示）。（质子的距离用 $10^{-15}\,\mathrm m$。）}
\answerarea
\problem{10.6}{算出 ${}^7_3\mathrm{Li}(\mathrm p,\alpha){}^4_2\mathrm{He}$ 的反应能。有关同位素的质量如下：${}^1_1\mathrm H$，1.007825；${}^4_2\mathrm{He}$，4.002603；${}^7_3\mathrm{Li}$，7.015999。}
\answerarea\clearpage
\problem{10.7}{在第 6 题的核反应中，如果以 $1\,\mathrm{MeV}$ 的质子打击 $\mathrm{Li}$，向正垂直于质子束的方向观测到 ${}^4_2\mathrm{He}$ 的能量有多大？}
\answerarea
\problem{10.8}{试计算 1 克 ${}^{235}\mathrm U$ 裂变时全部释放的能量约等于多少煤在空气中燃烧所放出的热能（煤的燃烧热约等于 $33\times10^6$ 焦耳/千克；$1\,\mathrm{MeV}=1.6\times10^{-13}$ 焦耳）。}
\answerarea\clearpage
\problem{10.9}{计算按照（10.8-1）式中前四式的核聚变过程用去 1 克氘所放出的能量约等于多少煤在空气中燃烧所放出的热能（煤的燃烧热同第 8 题）。
原书式（10.8-1）的前四式为：
\[
\begin{aligned}
{}^2\mathrm H+{}^2\mathrm H&\longrightarrow{}^3\mathrm{He}+{}^1\mathrm n+3.25\,\mathrm{MeV},\\
{}^2\mathrm H+{}^2\mathrm H&\longrightarrow{}^3\mathrm H+{}^1\mathrm H+4.00\,\mathrm{MeV},\\
{}^3\mathrm H+{}^2\mathrm H&\longrightarrow{}^4\mathrm{He}+{}^1\mathrm n+17.6\,\mathrm{MeV},\\
{}^3\mathrm{He}+{}^2\mathrm H&\longrightarrow{}^4\mathrm{He}+{}^1\mathrm H+18.3\,\mathrm{MeV}.
\end{aligned}
\]}
\answerarea
\problem{10.10}{包围等离子体的磁通量密度 $B$ 是 $2\,\mathrm T$，算出被围等离子体的压强。}
\answerarea\clearpage
